\section{Related Work}
\label{sec:related}

\paragraph{Prompt Injection Attacks.}
Perez et al.~\cite{perez2022ignore} first showed LLMs cannot reliably separate instructions from data; Greshake et al.~\cite{greshake2023not} extended this to \emph{indirect} injection via external content. Tensor Trust~\cite{tensortrust2024} quantified injection robustness through an adversarial game. HijackRAG~\cite{hijackrag2025} demonstrates that retrieval-augmented systems are particularly vulnerable. ToolHijacker~\cite{toolhijacker2026} shows tool-selection manipulation attacks against multiple defenses. Benchmarks including InjecAgent~\cite{injecagent2024}, AgentDojo~\cite{agentdojo2024}, ASB~\cite{asb2025}, and CyberSecEval~\cite{cyberseceval2024} provide standardized evaluation, and OWASP~\cite{owasp-llm2025} ranks prompt injection as the top LLM risk.

\paragraph{Filtering, Guardrails, and Training-Time Defenses.}
Pattern filters (Rebuff~\cite{rebuff-defense}, LLM Guard~\cite{llm-guard}, PromptGuard~\cite{prompt-guard}) and prompt-level guardrails (NeMo~\cite{nemo-guardrails}, LangChain~\cite{langchain-guardrails}) achieve moderate detection but cannot distinguish trusted from untrusted arguments without provenance. Training-time interventions (StruQ~\cite{struq2024}, SecAlign~\cite{secalign2024}, Jatmo~\cite{jatmo2024}), instruction hierarchy~\cite{instruction-hierarchy2024}, and MELON~\cite{melon2025} (masked re-execution) are model-internal and cannot protect already-deployed models. IsolateGPT~\cite{isolategpt2025} provides execution isolation at the system level. AgentFuzz~\cite{agentfuzz2025} uses taint-style analysis for vulnerability \emph{detection} but not runtime enforcement. Spotlighting~\cite{spotlighting2024} reduces injection via prompt delimiters but provides no tool-call enforcement.

\paragraph{Provenance-Aware Tool-Call Enforcement.}
The most relevant concurrent work shares \sys{}'s core insight: conditioning tool-call authorization on data lineage. We compare these systems honestly below.

\emph{CaMeL}~\cite{camel2025} (Google DeepMind, March 2025) maintains a full data-flow graph where each variable carries capability metadata specifying origin and access rights; a custom Python interpreter recursively traverses the entire dependency graph before every tool call. In STRICT mode, even implicit information flows through conditionals are tracked. Enforcement is deterministic: pure Python policy functions return \texttt{Allowed()} or \texttt{Denied(reason)} with no LLM in the loop, achieving 0\% ASR on AgentDojo~\cite{agentdojo2024}. CaMeL's \emph{structural isolation} (dual-LLM architecture with Privileged and Quarantined LLMs) prevents untrusted data from reaching the planner, making paraphrase resilience unnecessary. The tradeoff is reduced utility (48--77\% task completion~\cite{opcamel2025,agentarmor2025}) and a requirement for custom interpreter infrastructure.

\emph{FIDES}~\cite{fides2025} (Microsoft Research, May 2025) implements formal information-flow control with dynamic taint-tracking of integrity and confidentiality labels through a lattice model---remarkably similar to \sys{}'s two-element trust lattice. FIDES provides formal noninterference proofs and achieves 0\% policy-violating injections on AgentDojo. Like CaMeL, FIDES relies on structural separation of trusted and untrusted data flows.

\emph{PCAS}~\cite{pcas2026} (February 2026) models the system state as a dependency graph capturing causal relationships among tool calls, results, and messages. A Datalog-derived declarative policy language and reference monitor intercept all actions before execution---architecturally closest to \sys{}'s YAML-policy-plus-sidecar approach. PCAS tracks transitive information flow but does not address LLM paraphrasing.

\emph{ARM}~\cite{arm2026} (April 2026) implements a provenance graph with field-level tracking, an integrity lattice, and sidecar proxy enforcement. ARM extends the paradigm by modeling denied actions as first-class provenance events with counterfactual edges, addressing ``causality laundering'' attacks where adversaries chain denied-then-retried operations.

\emph{Progent}~\cite{progent2025} (UC Berkeley, April 2025) provides a JSON-based DSL for privilege control with deterministic \texttt{SecureCall} wrappers at every tool invocation but does not track data provenance.

Noureddine et al.~\cite{designpatterns2025} codify provenance-aware execution as one of six standard \emph{design patterns} for securing LLM agents, positioning the approach as established practice. PFI~\cite{pfi2025} separates trusted/untrusted agents with data proxies. Agent-Sentry~\cite{agentsentry2026} constructs functionality graphs from execution provenance for behavioral bounding. AgentArmor~\cite{agentarmor2025} achieves 72\% utility (vs.\ CaMeL's 48\%) with 3\% ASR.

\paragraph{Data Provenance and Taint Tracking.}
Why-provenance~\cite{buneman2001provenance}, the W3C PROV-DM standard~\cite{prov-dm}, dynamic taint analysis~\cite{taint-analysis-survey,taintdroid}, and DIFC systems~\cite{difc-survey,flume} provide foundational models. PROV-AGENT~\cite{provagent2025} applies W3C PROV-DM to agentic workflows but for \emph{observability} only---not security enforcement. The ``Crossing NL/PL'' analysis~\cite{nlpldivide2026} characterizes information flow through LLM calls including paraphrasing, using taxonomy-based filtering rather than embeddings. DIPPER~\cite{dipper2023} demonstrates embedding-based retrieval to trace text through paraphrasing for AI-text detection---the closest technique to \sys{}'s resolution pipeline, applied in an adjacent domain.

\paragraph{Capability-Based and IoT Security.}
\sys{} adapts capability-based security~\cite{miller2006capability,capsicum} to LLM tool calls and extends context-aware IoT access control (ContextIoT~\cite{jia2017contexiot}, IoTGuard~\cite{iotguard}) with provenance tracking.

\paragraph{Positioning \sys{}.}
Table~\ref{tab:feature-comparison} compares \sys{} with provenance-aware defenses. The distinguishing characteristics of \sys{} are: (1)~\emph{embedding-based provenance resolution}---no other security system uses sentence-embedding similarity for paraphrase-resilient provenance tracking, enabling enforcement without structural isolation; (2)~\emph{sidecar proxy deployment}---no model modification, no custom interpreter, framework-agnostic (shared with ARM but distinct from CaMeL, FIDES, PCAS); (3)~\emph{W3C PROV-DM grounding}---the only enforcement system using a standardized provenance model (PROV-AGENT uses W3C PROV for observability only). Systems based on structural isolation (CaMeL, FIDES) achieve stronger worst-case guarantees but sacrifice utility; \sys{} targets the complementary setting where strict isolation is impractical.

\begin{table}[t]
\centering
\caption{Comparison with provenance-aware defenses. \textbf{Prov.}: provenance graph/DAG; \textbf{Trust}: integrity/trust propagation; \textbf{Det.}: deterministic enforcement; \textbf{Sem.}: paraphrase-resilient resolution; \textbf{Sidecar}: framework-agnostic external proxy; \textbf{Formal}: standards-based or formally verified model.}
\label{tab:feature-comparison}
\scriptsize
\resizebox{\columnwidth}{!}{%
\begin{tabular}{lcccccc}
\toprule
\textbf{System} & \textbf{Prov.} & \textbf{Trust} & \textbf{Det.} & \textbf{Sem.} & \textbf{Sidecar} & \textbf{Formal} \\
\midrule
CaMeL~\cite{camel2025}         & \cmark & \cmark & \cmark & \xmark & \xmark & \xmark \\
FIDES~\cite{fides2025}          & \cmark & \cmark & \cmark & \xmark & \xmark & \cmark$^\dag$ \\
PCAS~\cite{pcas2026}            & \cmark & \cmark & \cmark & \xmark & \xmark & \xmark \\
ARM~\cite{arm2026}              & \cmark & \cmark & \cmark & \xmark & \cmark & \xmark \\
Progent~\cite{progent2025}      & \xmark & \xmark & \cmark & \xmark & \xmark & \xmark \\
Invariant~\cite{invariant2024}  & \xmark & \xmark & \cmark & \xmark & \xmark & \xmark \\
\midrule
\textbf{\sys{} (Ours)}         & \cmark & \cmark & \cmark & \cmark & \cmark & \cmark$^\ddag$ \\
\bottomrule
\multicolumn{7}{l}{\scriptsize $^\dag$Noninterference proofs. \quad $^\ddag$W3C PROV-DM standard.}
\end{tabular}%
}
\end{table}
